% ============================================
% LUC-18 — O Que Ronaldo Não Disse (edição revista R1075)
% ============================================
\fragtitle{luc}{L}{LUC}{18}{O Que Ronaldo Não Disse}
\markboth{LUC-18 --- O Que Ronaldo Não Disse}{}

\noindent \textit{Arquivo Público Mineiro --- 14 de abril de 2026, 19h05}

\noindent A sala de Ronaldo ficava no fim do corredor, depois da mapoteca. Lúcia conhecia o caminho de cor. Havia uma janela alta, com vista para o pátio interno, e uma estante de madeira escura que ele nunca terminara de organizar. Sobre a mesa, o mesmo abajur de latão, o mesmo porta-canetas de vidro, a mesma fotografia da mãe. Nada mudara em oito anos, exceto o quadro na parede: uma reprodução de uma vala comum, tirada de um livro de fotografia.

\noindent ``Você mandou meu nome para a comissão de ética'', ele disse, sem levantar os olhos.

\noindent ``Mandei.''

\noindent ``Por quê?''

\noindent ``Porque você me deu acesso a um documento que não estava catalogado. Porque você sabia que ele estava lá e não me contou. Porque você está com medo, Ronaldo. E medo não é motivo para colocar outra pessoa na frente.''

\noindent Ele levantou os olhos. Estava cansado. Havia uma olheira funda, nova, que não existia em janeiro. O olho direito, sob a luz do abajur, parecia mais claro. Lúcia disse a si mesma que era impressão.

\noindent ``Meu pai foi internado no Colônia em 1968'', ele disse.

\noindent Lúcia ficou parada.

\noindent ``Ele tinha 34 anos. Diagnóstico: esquizofrenia. Ficou lá seis meses. Nunca falou sobre isso. Quando eu era criança, ele tinha crises. Ficava horas no quintal, olhando para o muro. Minha mãe dizia que era o sol. Uma vez, ele me pegou pelo braço e disse: \textit{Não deixa ninguém ler o que está no meu armário.} Eu tinha nove anos. Não entendi. Depois esqueci.''

\noindent ``O que tinha no armário?''

\noindent ``Um caderno. Capa preta. Ele queimou antes de morrer. Eu vi. Ele colocou no quintal, jogou álcool e acendeu. Ficou olhando até virar cinza. Depois me olhou e disse: \textit{Agora está limpo.}''

\noindent Lúcia sentiu o cheiro de terra de novo. Não no ar. Na memória. Um cheiro que ela nunca tinha sentido na infância, mas que agora parecia ter estado sempre ali, escondido atrás das coisas.

\noindent ``Por que você me mandou o e-mail?''

\noindent Ronaldo abriu a gaveta. Tirou um envelope pardo, gasto, com o carimbo do arquivo. Colocou sobre a mesa.

\noindent ``Porque eu li a primeira página'', ele disse. ``Em 2019. Quando o diário chegou aqui, na primeira devolução. Eu abri. Só a primeira página. Vi a letra. Vi o nome. Fechei. Não contei a ninguém. Coloquei numa caixa e não abri mais. Mas o cheiro ficou comigo. Toda vez que eu passava perto da caixa, sentia. Toda vez que eu tentava dormir, sonhava com uma vala. Eu não sabia o que era. Agora sei.''

\noindent ``O que você sabe?''

\noindent ``Que o diário escolhe. E que eu não podia ser o próximo. Eu tenho uma filha, Lúcia. Eu não podia.'' Ele fez uma pausa. ``Então eu procurei alguém que pudesse. Alguém que soubesse ler o que estava ali. Alguém que não fosse quebrar.''

\noindent Lúcia sentiu o rosto esquentar. ``Você me usou.''

\noindent ``Eu te escolhi.''

\noindent ``É a mesma coisa.''

\noindent ``Não é.'' Ele levantou a voz, depois baixou. ``Eu te escolhi porque você era a única pessoa que eu conhecia que não ia rir. Que ia entender. Que ia fazer as perguntas certas. Eu não te escolhi para morrer. Eu te escolhi para ler.''

\noindent ``E o que acontece depois que eu leio?''

\noindent Ronaldo não respondeu. Olhou para a fotografia da mãe. A luz do abajur tremeluziu.

\noindent ``Você leu o diário inteiro?'' ela perguntou.

\noindent ``Não.''

\noindent ``Nunca?''

\noindent ``Nunca.''

\noindent ``Então como sabe que ele escolhe?''

\noindent Ele abriu a boca para responder. Não saiu som. Ficou olhando para ela, com uma expressão que Lúcia não conseguiu decifrar: medo, culpa, ou algo mais antigo. O olho direito dele brilhou. Um brilho amarelo, rápido, como um fósforo aceso e apagado.

\noindent Lúcia se levantou. Pegou o envelope.

\noindent ``Eu vou ler'', ela disse. ``Mas não porque você me escolheu. Porque eu quero saber.''

\noindent Na porta, Ronaldo falou:

\noindent ``Lúcia.''

\noindent Ela parou.

\noindent ``Se você encontrar meu pai lá dentro'', ele disse, ``não fala com ele. Ele não gosta de ser encontrado.''

\noindent Lúcia saiu. No corredor, o cheiro de terra era tão forte que ela precisou se apoiar na parede. O envelope em sua mão estava quente. Quando chegou ao carro, abriu. O envelope não trazia o diário. Trazia a carta de 1853. E, no verso, numa letra que ela conhecia, havia uma anotação:

\noindent ``Agora você sabe por que eu não podia ler.''

> **↪ Rotas:** **LUC-17** · **LUC-19** · **DOC-27**
